\textbf{Phần 1. Trắc nghiệm lựa chọn}
\setcounterpageref{ex}{0}

\begin{ex} %[1D1V06]
Tìm số lớn nhất có một chữ số.
\choice
{$8$}
{\True $9$}
{$10$}
{$7$}
\loigiai{
Trong các số tự nhiên có một chữ số từ 0 đến 9, số lớn nhất là số $9$.
\textbf{Đáp số:} $9$.\par
\textbf{Đáp án của câu hỏi là B}
}
\end{ex}

\begin{ex} %[1D1V06]
Tính tổng của số bé nhất có hai chữ số và số bé nhất có một chữ số.
\choice
{\True $10$}
{$0$}
{$20$}
{$9$}
\loigiai{
Số bé nhất có hai chữ số là $10$.\nSố bé nhất có một chữ số là $0$.\nTổng của hai số là: $10 + 0 = 10$.
\textbf{Đáp số:} $10$.\par
\textbf{Đáp án của câu hỏi là A}
}
\end{ex}

\begin{ex} %[1D1V06]
Tìm hiệu giữa số lớn nhất có một chữ số và số bé nhất khác 0.
\choice
{$7$}
{$9$}
{\True $8$}
{$6$}
\loigiai{
Số lớn nhất có một chữ số là $9$.\nSố bé nhất khác 0 là số $1$.\nHiệu của hai số là: $9 - 1 = 8$.
\textbf{Đáp số:} $8$.\par
\textbf{Đáp án của câu hỏi là C}
}
\end{ex}

\begin{ex} %[1D1V06]
Bạn An có 9 viên kẹo, An cho bạn Bình 3 viên và cho bạn Cúc 2 viên. Hỏi An còn lại bao nhiêu viên kẹo?
\choice
{$3$}
{$5$}
{$2$}
{\True $4$}
\loigiai{
Tổng số kẹo An đã cho hai bạn là: $3 + 2 = 5$ (viên kẹo).\nSố kẹo An còn lại là: $9 - 5 = 4$ (viên kẹo).
\textbf{Đáp số:} $4$.\par
\textbf{Đáp án của câu hỏi là D}
}
\end{ex}

\begin{ex} %[1D1V06]
Trong bãi đỗ xe có 10 chiếc xe gồm xe máy và xe đạp. Biết số xe máy là 7 chiếc. Hỏi số xe máy nhiều hơn số xe đạp bao nhiêu chiếc?
\choice
{$3$}
{\True $4$}
{$5$}
{$2$}
\loigiai{
Số xe đạp trong bãi là: $10 - 7 = 3$ (chiếc xe đạp).\nSố xe máy nhiều hơn số xe đạp là: $7 - 3 = 4$ (chiếc).
\textbf{Đáp số:} $4$.\par
\textbf{Đáp án của câu hỏi là B}
}
\end{ex}

\begin{ex} %[1D1V06]
Có bao nhiêu số tự nhiên lớn hơn 3 và bé hơn 10?
\choice
{\True $6$}
{$5$}
{$7$}
{$4$}
\loigiai{
Các số đó là: $4, 5, 6, 7, 8, 9$.\nĐếm các số trên có tất cả $6$ số.
\textbf{Đáp số:} $6$.\par
\textbf{Đáp án của câu hỏi là A}
}
\end{ex}

\begin{ex} %[1D1V06]
Cho dãy số cách đều: $0, 2, 4, 6, \dots, 10$. Số thích hợp điền vào dấu ba chấm là số nào?
\choice
{$7$}
{$9$}
{\True $8$}
{$6$}
\loigiai{
Quy luật dãy số tăng dần cách đều 2 đơn vị: $0, 2, 4, 6, 8, 10$.\nSố còn thiếu là số $8$.
\textbf{Đáp số:} $8$.\par
\textbf{Đáp án của câu hỏi là C}
}
\end{ex}

\begin{ex} %[1D1V06]
Cho dãy số đếm lùi cách đều 2 đơn vị: $9, 7, 5, \dots, 1$. Số thích hợp điền vào dấu ba chấm là số nào?
\choice
{$2$}
{$4$}
{$1$}
{\True $3$}
\loigiai{
Mỗi số đứng sau kém số đứng trước 2 đơn vị: $5 - 2 = 3$ và $3 - 2 = 1$.\nSố cần điền là số $3$.
\textbf{Đáp số:} $3$.\par
\textbf{Đáp án của câu hỏi là D}
}
\end{ex}

\begin{ex} %[1D1V06]
Tìm một số tự nhiên biết rằng khi lấy số đó cộng với 4 rồi trừ đi 2 thì được kết quả bằng 7.
\choice
{$4$}
{\True $5$}
{$6$}
{$3$}
\loigiai{
Lấy kết quả là 7 cộng ngược lại 2: $7 + 2 = 9$.\nSau đó lấy 9 trừ đi 4: $9 - 4 = 5$.\nThử lại: $5 + 4 - 2 = 7$ (đúng). Vậy số đó là $5$.
\textbf{Đáp số:} $5$.\par
\textbf{Đáp án của câu hỏi là B}
}
\end{ex}

\begin{ex} %[1D1V06]
Có 10 que tính, bạn Mai xếp được 2 hình vuông riêng biệt (mỗi hình vuông dùng 4 que tính). Hỏi Mai còn thừa bao nhiêu que tính?
\choice
{\True $2$}
{$1$}
{$3$}
{$4$}
\loigiai{
Số que tính dùng để xếp 2 hình vuông là: $4 + 4 = 8$ (que tính).\nSố que tính còn thừa là: $10 - 8 = 2$ (que tính).
\textbf{Đáp số:} $2$.\par
\textbf{Đáp án của câu hỏi là A}
}
\end{ex}

\begin{ex} %[1D1V06]
Nhà bạn Linh nuôi 4 con thỏ trắng và số thỏ xám nhiều hơn thỏ trắng 2 con. Hỏi nhà bạn Linh nuôi tất cả bao nhiêu con thỏ?
\choice
{$0$}
{$20$}
{\True $10$}
{$9$}
\loigiai{
Số thỏ xám là: $4 + 2 = 6$ (con thỏ).\nTổng số thỏ nhà Linh nuôi là: $4 + 6 = 10$ (con thỏ).
\textbf{Đáp số:} $10$.\par
\textbf{Đáp án của câu hỏi là C}
}
\end{ex}

\begin{ex} %[1D1V06]
Hôm nay là Thứ Ba ngày 5. Hỏi 3 ngày sau là ngày mùng mấy?
\choice
{$7$}
{$9$}
{$6$}
{\True $8$}
\loigiai{
Ngày sau đó 3 ngày là ngày: $5 + 3 = 8$ (ngày mùng 8).
\textbf{Đáp số:} $8$.\par
\textbf{Đáp án của câu hỏi là D}
}
\end{ex}

\textbf{Phần 2. Trắc nghiệm đúng sai}
\setcounterpageref{ex}{0}

\begin{ex} %[1D1V06]
Xét tính đúng hoặc sai của các khẳng định toán học sau đây:
\choiceTF
{Kết quả của bài toán «Tìm số tự nhiên $x$ thỏa mãn: $8 - x = 2 + 3$. Giá trị của $x$ là bao nhiêu» là $2$.}
{\True Kết quả của bài toán «Tìm số tự nhiên $y$ thỏa mãn: $y + 4 = 10 - 1$. Giá trị của $y$ là bao nhiêu» là $5$.}
{\True Một cây nến đang cháy có chiều dài 10 cm. Sau một giờ nến bị cháy ngắn bớt đi 3 cm. cây nến còn lại dài $7\text{ cm}$ xăng-ti-mét.}
{Trong giỏ có 8 quả lê. Mẹ chia đều cho hai chị em mỗi người một nửa số lê. mỗi người được $5$ quả lê.}
\loigiai{
\begin{itemize}
    \item Ý a) Sai: Kết quả đúng là $3$, không phải $2$. Tính vế phải: $2 + 3 = 5$.\nTa có: $8 - x = 5 \Rightarrow x = 8 - 5 = 3$.
\textbf{Đáp số:} $3$.
    \item Ý b) Đúng: Tính vế phải: $10 - 1 = 9$.\nTa có: $y + 4 = 9 \Rightarrow y = 9 - 4 = 5$.
\textbf{Đáp số:} $5$.
    \item Ý c) Đúng: Chiều dài còn lại của cây nến là: $10 - 3 = 7$ (cm).
\textbf{Đáp số:} $7$.
    \item Ý d) Sai: Kết quả đúng là $4$, không phải $5$. Vì hai chị em nhận số lê bằng nhau và $4 + 4 = 8$ nên mỗi người được $4$ quả lê.
\textbf{Đáp số:} $4$.
\end{itemize}
\begin{itemize}
    \item a) \textbf{Sai}.
    \item b) \textbf{Đúng}.
    \item c) \textbf{Đúng}.
    \item d) \textbf{Sai}.
\end{itemize}
}
\end{ex}

\begin{ex} %[1D1V06]
Xét tính đúng hoặc sai của các bài toán và phép tính sau đây:
\choiceTF
{\True Một nhóm có 7 bạn đi dã ngoại. Có 4 bạn đội mũ đỏ, số bạn còn lại đội mũ xanh. số bạn đội mũ đỏ nhiều hơn số bạn đội mũ xanh $1$ bạn.}
{Số liền trước của số lớn nhất có một chữ số là số $9$.}
{\True Số liền sau của số bé nhất có hai chữ số là số $11$.}
{Một cây nến đang cháy có chiều dài 10 cm. Sau một giờ nến bị cháy ngắn bớt đi 3 cm. cây nến còn lại dài $6\text{ cm}$ xăng-ti-mét.}
\loigiai{
\begin{itemize}
    \item Ý a) Đúng: Số bạn đội mũ xanh là: $7 - 4 = 3$ (bạn).\nSố bạn đội mũ đỏ nhiều hơn số bạn đội mũ xanh là: $4 - 3 = 1$ (bạn).
\textbf{Đáp số:} $1$.
    \item Ý b) Sai: Kết quả đúng là $8$, không phải $9$. Số lớn nhất có một chữ số là số $9$.\nSố đứng liền trước số 9 là: $9 - 1 = 8$.
\textbf{Đáp số:} $8$.
    \item Ý c) Đúng: Số bé nhất có hai chữ số là $10$.\nSố liền sau của 10 là: $10 + 1 = 11$.
\textbf{Đáp số:} $11$.
    \item Ý d) Sai: Kết quả đúng là $7$, không phải $6$. Chiều dài còn lại của cây nến là: $10 - 3 = 7$ (cm).
\textbf{Đáp số:} $7$.
\end{itemize}
\begin{itemize}
    \item a) \textbf{Đúng}.
    \item b) \textbf{Sai}.
    \item c) \textbf{Đúng}.
    \item d) \textbf{Sai}.
\end{itemize}
}
\end{ex}

\textbf{Phần 3. Trả lời ngắn}
\setcounterpageref{ex}{0}

\begin{ex} %[1D1V06]
Cho ba số: $3, 5, 2$. Lấy số lớn nhất cộng với số bé nhất được kết quả bằng bao nhiêu?
\shortans{7}
\loigiai{
\textbf{Đáp số:} $7$\par
Trong ba số trên, số lớn nhất là $5$ và số bé nhất là $2$.\nTổng của hai số là: $5 + 2 = 7$.
\textbf{Đáp số:} $7$.
}
\end{ex}

\begin{ex} %[1D1V06]
Cho ba số: $10, 4, 7$. Lấy số lớn nhất trừ đi số bé nhất được kết quả bằng bao nhiêu?
\shortans{6}
\loigiai{
\textbf{Đáp số:} $6$\par
Trong ba số trên, số lớn nhất là $10$ và số bé nhất là $4$.\nHiệu của hai số là: $10 - 4 = 6$.
\textbf{Đáp số:} $6$.
}
\end{ex}

\begin{ex} %[1D1V06]
Có 9 quả bóng được chia thành hai phần sao cho phần thứ nhất có nhiều hơn phần thứ hai 1 quả. Hỏi phần thứ nhất có bao nhiêu quả bóng?
\shortans{5}
\loigiai{
\textbf{Đáp số:} $5$\par
Hai số có tổng bằng 9 và hiệu bằng 1 là $5$ và $4$ (vì $5 + 4 = 9$ và $5 - 4 = 1$).\nPhần thứ nhất có $5$ quả bóng.
\textbf{Đáp số:} $5$.
}
\end{ex}

\begin{ex} %[1D1V06]
Bình có 6 viên bi, An có nhiều hơn Bình 2 viên bi. Cả hai bạn gom bi lại và chia đều cho nhau. Hỏi mỗi bạn nhận được bao nhiêu viên bi?
\shortans{7}
\loigiai{
\textbf{Đáp số:} $7$\par
Số bi của An là: $6 + 2 = 8$ (viên bi).\nTổng số bi của cả hai bạn là: $6 + 8 = 14$ viên bi (chia đều mỗi bạn 7 viên).
\textbf{Đáp số:} $7$.
}
\end{ex}

\begin{ex} %[1D1V06]
Tìm số tự nhiên $a$ biết rằng: $a$ lớn hơn 6 và $a$ bé hơn 8.
\shortans{7}
\loigiai{
\textbf{Đáp số:} $7$\par
Số tự nhiên nằm giữa 6 và 8 là số $7$.
\textbf{Đáp số:} $7$.
}
\end{ex}

\begin{ex} %[1D1V06]
Một sợi dây dài 10 cm, bị cắt bớt đi một đoạn còn lại 6 cm. Hỏi đoạn dây bị cắt dài bao nhiêu xăng-ti-mét?
\shortans{4}
\loigiai{
\textbf{Đáp số:} $4$\par
Đoạn dây bị cắt dài là: $10 - 6 = 4$ (cm).
\textbf{Đáp số:} $4$.
}
\end{ex}
